\honest{When (Not) To Use Probabilities}{Eliezer Yudkowsky}{2008}

I do not always advocate making up probabilities. The laws of probability are laws, but often they are too hard to compute. \nb{That holds under the conditions the post states: exact probabilistic inference is NP-hard (Cooper, 1990).} To catch a ball, trust your brain's built-in mechanisms rather than a distribution over landing spots. The Dutch book arguments still apply, but catching the ball matters more than being consistent.

So I advise against making up probabilities without some decent basis. ``Numbers should come from numbers.'' If you turn a gut feeling into ``one in a thousand'', the event will not happen one time in a thousand. \nb{A study quoted in Yudkowsky's own chapter on global risks found answers given odds of 1000 to 1 right 81\% of the time.}

My example is an argument from the Global Catastrophic Risks conference. \nb{Almost certainly the talk ``Probing the Improbable'' by Hillerbrand, Sandberg and Ord, three days earlier.} The speaker put at least 1 in 1000 on the LHC safety papers being wrong, and at least 1 in 1000 on the world being destroyed if they were. I object to ``these numbers pulled out of thin air.'' \nb{The second figure was a guess, as the speakers' written paper admits. The first rested, by a journalist's report of the talk, on rates of retraction of journal articles, a number from numbers. Yudkowsky's own 2006 draft chapter had said that experts who state one in a million ``have undoubtedly been wrong more often than one time in a million.''}

I propose instead to debate the general rule of banning physics experiments that cannot be proved safe. \nb{Two days earlier Yudkowsky had reported that the proposer ``disagreed that their policy cashed out to a blanket ban on physics experiments.''} The speakers' wish to do more analysis and then switch the LHC on I call ``the most disingenuous part of the argument'', since no paper can remove the doubt. \nb{Their paper, three months later, grants that the doubt ``can never be completely removed'' and argues that independent arguments shrink it. That is a consistent position; the post gives no ground for a charge of bad faith.} The speakers would have made up the same numbers anyway, I say, and I give this a probability of 75\%.

I admit an inconsistency: I would be more alarmed by a lottery device with a one-in-a-million chance of destroying the world than by the LHC, yet I could not make a million statements as sure as ``The Large Hadron Collider will not destroy the world'' and be wrong only about once. \nb{By that test the author's chance of error is above about one in a million, close to the speaker's figures, as the post goes on to admit.} I keep the inconsistency, because the goal is to win.

That is why my posts say ``maybe'' and ``probably'' rather than ``40\%'' and ``70\%''. ``Think of how silly that would look.'' I think I would do worse with numbers. \nb{Sherman Kent of the CIA answered the same objection in 1964; colleagues who had agreed on ``serious possibility'' turned out to mean odds from 20 to 80 to 80 to 20.}
